\documentclass[11pt,a4paper]{article}

\usepackage[T1]{fontenc}
\usepackage[utf8]{inputenc}
\usepackage[ngerman]{babel}
\usepackage[a4paper,margin=2.5cm]{geometry}
\usepackage{xcolor}
\usepackage{graphicx}
\usepackage{booktabs}
\usepackage{hyperref}
\usepackage{minted}

\graphicspath{{../screenshots/}}

\title{ATS ULP V1 -- Homing Model Dokumentation}
\author{Florian}
\date{\today}

\begin{document}

\maketitle

\section{Projektstruktur}
Das Projekt ist als kleines Python-Paket organisiert. Die Kernlogik liegt in den Modulen unter \texttt{homing\_code/}. Der Einstiegspunkt liegt auf der Repository-Wurzelebene in \texttt{main.py}.

\section{Code-Referenzen}
\subsection{Datenmodelle}
\inputminted[fontsize=\small,linenos,breaklines]{python}{../../homing_code/models.py}

\subsection{Geometrie}
\inputminted[fontsize=\small,linenos,breaklines]{python}{../../homing_code/geometry.py}

\subsection{Pairing und Vektoren}
\inputminted[fontsize=\small,linenos,breaklines]{python}{../../homing_code/pairing.py}

\subsection{Visualisierung}
\inputminted[fontsize=\small,linenos,breaklines]{python}{../../homing_code/visualization.py}

\subsection{Demo-Start}
\inputminted[fontsize=\small,linenos,breaklines]{python}{../../main.py}

\section{Screenshots}
Alle Screenshots für die Dokumentation werden in \texttt{docs/screenshots/} abgelegt und von dort eingebunden.

\begin{figure}[h]
  \centering
  \fbox{\parbox{0.8\linewidth}{Screenshot-Platzhalter. Datei in \texttt{docs/screenshots/} ablegen und hier mit \texttt{\textbackslash includegraphics} referenzieren.}}
  \caption{Platzhalter für Visualisierungen aus dem Homing-Projekt.}
\end{figure}

\end{document}